\par
Para (5), $R$ es formalmente liso sobre $\Lambda$, y por 4.4 y 4.11
la sucesión
\[
0\longrightarrow H^1(X_0,\underline D^0_{X_0})
 \longrightarrow[\underline M_{X_0}(k[\varepsilon])]
 \longrightarrow H^0(X_0,\underline D^1_{X_0})
 \longrightarrow0
\]
es exacta. Así,
\[
\begin{aligned}
\dim\underline M_{X_0}
&=\dim_kH^1(X_0,\underline D^0_{X_0})
 +\dim_kH^0(X_0,\underline D^1_{X_0})\\
&=\dim_kH^1(X_0,\underline D^0_{X_0})+m.
\end{aligned}
\]
Basta, pues, mostrar que
\[
\dim_kH^1(X_0,\underline D^0_{X_0})+m
=3g-3+\dim_kH^0(X_0,\underline D^0_{X_0}),
\]
o equivalentemente
\[
\chi(\underline D^0_{X_0})=3-3g+m,
\]
donde, para un $\mathcal O_{X_0}$-módulo coherente $\mathcal E$,
ponemos
\[
\chi(\mathcal E)
=\dim_kH^0(X_0,\mathcal E)-\dim_kH^1(X_0,\mathcal E).
\]

La característica de Euler de $\underline D^0_{X_0}$ no cambia al
extender el cuerpo base a su clausura algebraica, de modo que podemos
suponer que $k$ es algebraicamente cerrado. Sea $z$ un punto no liso
de $X_0$, pongamos
\[
A=\widehat{\mathcal O}_{X_0,z},
\]
y sea $A'$ la normalización de $A$. Por hipótesis,
\[
A\simeq k[[x,y]]/(f),
\]
donde $f$ es una forma cuadrática no degenerada. Sobre el cuerpo
algebraicamente cerrado $k$ podemos tomar $f=xy$, de modo que
\[
A\simeq k[[x,y]]/(xy).
\]
Se sigue inmediatamente que
\[
\ell(A'/A)=1,
\qquad
\operatorname{Der}_k(A,\mathfrak m_A)
=\operatorname{Der}_k(A,A).
\]
Sea $C$ el conductor de $A$ en $A'$. Como $A$ es una intersección
completa,
\[
\ell(A/C)=\ell(A'/A)=1,
\]
y por tanto $C=\mathfrak m_A$. En consecuencia,
\[
A'=\{a'\in K\mid a'\mathfrak m_A\subset\mathfrak m_A\},
\]
donde $K$ es el anillo total de fracciones de $A$.

Afirmamos que
\[
\operatorname{Der}_k(A,A)
 \subset\operatorname{Der}_k(A',A')
\]
dentro de $\operatorname{Der}_k(K,K)$. En efecto, sea
$D\in\operatorname{Der}_k(A,A)$. Para $a'\in A'$ y
$x\in\mathfrak m_A$, se tiene $xa'\in\mathfrak m_A$, de modo que
$D(xa')\in\mathfrak m_A$ y $a'D(x)\in\mathfrak m_A$. La identidad
\[
D(xa')=a'D(x)+xD(a')
\]
da, por tanto, $xD(a')\in\mathfrak m_A$ para todo
$x\in\mathfrak m_A$, de donde $D(a')\in A'$.

%%%%%%%%%%  scan idx 141  |  folio 130  %%%%%%%%%%

Además, la sucesión exacta
\[
0\longrightarrow\operatorname{Der}_k(A,A)
 \longrightarrow\operatorname{Der}_k(A',A')
 \longrightarrow A'/\mathfrak m_AA'\longrightarrow0
\]
puede escribirse
\[
0\longrightarrow\operatorname{Der}_k(A,A)
 \longrightarrow\operatorname{Der}_k(A',A')
 \longrightarrow A'/C\longrightarrow0.
\]
Si $X_0'$ es la normalización de $X_0$, estas sucesiones locales dan
una sucesión exacta
\[
0\longrightarrow\underline D^0_{X_0}
 \longrightarrow\underline D^0_{X_0'}
 \longrightarrow\mathcal O_{X_0'}/C\longrightarrow0.
\]
Como $X_0$ es localmente una intersección completa,
\[
\ell(\mathcal O_{X_0'}/C)
=2\ell(\mathcal O_{X_0'}/\mathcal O_{X_0})=2m,
\]
mientras que el teorema de Riemann--Roch da
\[
\chi(\underline D^0_{X_0'})=3-3(g-m).
\]
Por consiguiente,
\[
\begin{aligned}
\chi(\underline D^0_{\sisi{X^0_0}{X_0}})
&=\chi(\underline D^0_{X_0'})
 -\ell(\mathcal O_{X_0'}/C)\\
&=3-3(g-m)-2m=3-3g+m.
\end{aligned}
\]
Esto concluye la demostración.
\hfill$\square$

\par\medskip
\noindent\textbf{Observación 6.5.}
En la fórmula del Teorema 6.4(5), el número
$a=\dim_kH^0(X_0,\underline D^0_{X_0})$ se calcula fácilmente por el
teorema de Riemann--Roch:
\[
a=
\begin{cases}
0,&g\geq2,\\
1,&g=1,\\
3,&g=0.
\end{cases}
\]

Terminamos con una observación sobre la característica $2$. Si
$\operatorname{char}k=2$, no hay ninguna forma cuadrática no
degenerada en un número impar de variables. Equivalentemente, no hay
ninguna singularidad cuadrática no degenerada de dimensión par en
característica $2$. En este caso se puede considerar, en su lugar, una
forma cuadrática de no degeneración maximal.
